\documentclass{article}
\begin{document}

\begin{abstract}
Retrieval-augmented generation grounds a language model in retrieved evidence.
We report a system that retrieves passages from a domain corpus and conditions
generation on them, and we find that it reduces unsupported claims relative to
an unretrieved baseline. Our evaluation uses a single held-out set of 400
questions and reports exact-match accuracy only.
\end{abstract}

\section{Introduction}
Language models produce fluent text that is not always supported by any source.
Retrieval-augmented generation addresses this by conditioning generation on
passages fetched at inference time, so that a claim can in principle be traced
to a document. Prior work establishes the architecture; what remains unclear is
how much of the observed gain comes from retrieval itself rather than from the
longer context it supplies. We hold context length fixed and vary only whether
the passages are retrieved or sampled at random. The remainder of this paper
describes the corpus, the retriever, and the evaluation protocol.

\end{document}
